\section{One causal programme for a prepared family}\label{sec:uniform}
The quantifier on the programme is essential when an experimental parameter is inaccessible. A separate feasible allocation at each parameter need not be implementable by one observer. We first give an exact common-programme criterion and then an effective, two-sided value theorem. The parameter is fixed before preparation and remains fixed throughout all calls. In particular, taking a worst case anew at each call is a different experiment.

\subsection{The common parameter and phase constraint}
Let $\Theta$ be a nonempty parameter set. At each $\theta$, assume the prepared compact predictive experiment of Section~\ref{sec:experiment}, with carrier $K_{\theta,t}$, preparation $q_{\theta,0}$, instrument $(J_{\theta,t},F_{\theta,t})$, and terminal loss $\ell_\theta$. The finite action sets, action-indexed report spaces, horizon $n$ and decision space $U$ are common. The modeller knows this family; the observer does not receive $\theta$. Assume, in addition, that for each action there is one finite Borel measure $\sigma_a$ dominating every $J_{\theta,t}(q,a)$ at a legal phase and state. This joint domination is an additional hypothesis, not a consequence of parameterwise standard-Borel realization. The effective classes below have it. Normalizing a nonzero $\sigma_a$ is harmless.

Write $\mathcal G_{\aut}^{\Theta}(n,M)$ for the tuples of Definition~\ref{def:machines} which are legal and stop at $n$ under every $\theta$. All their components, including their terminal readouts, are parameter independent. Define
\begin{equation}\label{eq:robust_risk}
 R^\Theta_{\aut}(n,M)=\inf_{\gamma\in\mathcal G_{\aut}^{\Theta}(n,M)}
                         \sup_{\theta\in\Theta} R_\theta(\gamma).
\end{equation}
An empty class has value $+\infty$. The analogous external class allows phase-indexed rows, but not parameter-indexed rows. Terminal readouts are deterministic in the declared space $U$. If terminal lotteries are allowed, they are included in $U$, with their expected loss and sampling cost; no purification across parameters is assumed.

Fix $Z,z_0,D$, one action row $\alpha_z$, and one readout $u_j$ for each $j\in D$. A proposed family of occupancies consists of $p_{\theta,t,z}$ and $q_{\theta,t,z}\in K_{\theta,t}$ satisfying the initial, probability, exact-stopping and legality constraints \eqref{eq:occupancy} for every $\theta$. Thus a used action at $z$ is legal at \emph{every} occupied parameter--phase pair. Set
\[
 h_{\theta,t,z,a}
 =\frac{d[p_{\theta,t,z}\alpha_z(a)J_{\theta,t}(q_{\theta,t,z},a)]}{d\sigma_a}.
\]
These are nonnegative integrable functions. For a finite set
$\mathcal I\subset\Theta\times\{0,\ldots,n-1\}$ and
$\phi_{\theta,t,j}\in\aff(K_{\theta,t+1})$, consider
\begin{equation}\label{eq:uniform_dual}
\begin{split}
 &\sum_{(\theta,t)\in\mathcal I}\sum_j
       p_{\theta,t+1,j}\phi_{\theta,t,j}(q_{\theta,t+1,j})\\
 &\quad\le\sum_{z\notin D,a}\int
       \max_j\left\{\sum_{(\theta,t)\in\mathcal I}
        h_{\theta,t,z,a}(y)
        \phi_{\theta,t,j}(F_{\theta,t}(q_{\theta,t,z},a,y))\right\}\,\sigma_a(dy).
\end{split}
\end{equation}
The maximum is outside \emph{both} the parameter sum and the phase sum. For external widths, use fixed sets $Z_t$, common rows $\alpha_{t,z}$ and $k_t$, and require at each fixed $t$, for every finite $F\subset\Theta$,
\begin{equation}\label{eq:uniform_ext_dual}
 \sum_{\theta\in F,j}p_{\theta,t+1,j}\phi_{\theta,j}(q_{\theta,t+1,j})
 \le\sum_{z,a}\int\max_j\left\{\sum_{\theta\in F}
 h_{\theta,t,z,a}(y)\phi_{\theta,j}(F_{\theta,t}(q_{\theta,t,z},a,y))\right\}\,\sigma_a(dy).
\end{equation}
In this external display $h$ uses $\alpha_{t,z}$ in place of $\alpha_z$, and the same terminal readout is used across parameters. Let
$\mathfrak F^\Theta_{n,M}$ denote the autonomous proposals satisfying \eqref{eq:uniform_dual} and the stated common-component constraints. Define the actual-law invariant
\begin{equation}\label{eq:uniform_phi}
 \Phi^\Theta(n,M)=\inf_{\mathfrak F^\Theta_{n,M}}
       \sup_\theta\sum_{j\in D}p_{\theta,n,j}
                         \ell_\theta(q_{\theta,n,j},u_j).
\end{equation}
This is a risk-valued invariant. It does not separately minimize an oracle baseline or the terminal compression term. In particular, $u_j$ cannot be replaced by a different Bayes action at each parameter. The conditional predictive states are certificate variables under each actual law, not information furnished to the observer.

\begin{theorem}[Parameter-uniform causal acquired-geometry transfer]\label{thm:uniform}
For the dominated prepared family above, a proposed occupancy family is realized by one parameter-independent autonomous observer if and only if it belongs to $\mathfrak F^\Theta_{n,M}$. Any failed proposal has a strict witness \eqref{eq:uniform_dual} involving finitely many parameter--phase pairs and affine tests. Consequently
\begin{equation}\label{eq:uniform_exact}
                  R^\Theta_{\aut}(n,M)=\Phi^\Theta(n,M).
\end{equation}
For external widths the same statement holds with a common row at every fixed phase; the maximum then remains outside the parameter sum at that phase. No compactness of the parameter set is required for this exact assertion. Attainment is not asserted on arbitrary compact predictive carriers. If the family instead has a common finite-dimensional positive-instrument realization with joint domination and continuous compact terminal readouts, the optimum is attained and
\begin{equation}\label{eq:finite_model_value}
 R^\Theta_{\aut}(n,M)=
 \sup_{\varnothing\ne F\subset\Theta,\,|F|<\infty}R^F_{\aut}(n,M).
\end{equation}
An empty common class or a strict suboptimal risk target is then excluded already by finitely many fixed models. This conclusion needs neither continuity in the parameter nor normalized-filter stability. Theorem~\ref{thm:uniform_compactness} specifies and proves this positive-instrument clause.

Under the additional uniformly effective presentation and parameter moduli of Definition~\ref{def:uniform_effective}, this invariant has the following finite, globally valid certificates at the \emph{same} persistent budget. Choose a report partition, dyadic denominator $Q$, terminal loss net, a finite $s$-net $\Theta_s\subset\Theta$, and evaluation tolerance $\eta>0$. Enumerate the admissible programme set $\mathscr G_{h,Q,\rho}$ of Theorem~\ref{thm:effective}, using its common support graph. For each $\gamma$ and $\theta_i\in\Theta_s$, obtain
\[
 l_{\gamma,i}\le R_{\theta_i}(\gamma)\le u_{\gamma,i},\qquad u_{\gamma,i}-l_{\gamma,i}\le\eta,
 \quad L=\min_\gamma\max_i l_{\gamma,i},\quad U=\min_\gamma\max_i u_{\gamma,i}.
\]
If the class is nonempty, then
\begin{equation}\label{eq:uniform_interval}
 \boxed{\quad L-E_h-\beta_Q\ \le\ \Phi^\Theta(n,M)\ \le\ U+\Omega_n(s),
       \qquad U-L\le\eta.\quad}
\end{equation}
A programme minimizing $\max_i u_{\gamma,i}$ works on the original family without learning $\theta$ for free. Its worst-case excess over the optimal $M$-label Borel programme is at most
\begin{equation}\label{eq:uniform_gap}
                  E_h+\beta_Q+\Omega_n(s)+\eta.
\end{equation}
The list has at most $M2^M|U_\rho|^M(Q+1)^P$ members, with $P$ as in \eqref{eq:effective_constants}; each needs at most $|\Theta_s|$ certified risk evaluations. These certificates converge effectively from both sides. They do not merely bound the best programme in a selected finite list. Empty-class detection is exact under the support hypotheses. The corresponding external assertion uses \eqref{eq:effective_external} and its error budget.
\end{theorem}

The geometric content of \eqref{eq:uniform_exact} is the simultaneous realization of \emph{actually acquired} conditional laws with one finite programme. The quantitative content of \eqref{eq:uniform_interval} is a modulus-controlled converse to finite implementation. Neither replaces the acquired law by the ambient parameter dimension. The report net controls approximation of the physical interface, whereas the parameter net certifies a worst case over inaccessible models. Their sizes are computational costs, not substitutes for the risk geometry. The singleton-parameter theory and its matching regular and singular estimates are proved in Sections~\ref{sec:completion}--\ref{sec:singular}.

\subsection{Proof of the exact common-programme statement}
For fixed proposed occupancies, let $\mathscr K$ be the product, over $z\notin D$ and $a$, of the simplex-valued unit balls
\[
 \{(k_j)_j\in L^\infty(\sigma_a)^Z:k_j\ge0,\ \sum_j k_j=1\}.
\]
It is convex and weak-star compact: positivity and the row sum are closed inequalities and equalities against $L^1(\sigma_a)$. Map $\mathscr K$ to the product of all parameter--phase mass and affine-moment coordinates by
\begin{equation}\label{eq:uniform_moment}
 k\longmapsto\left(
   \sum_{z\notin D,a}\int k(j\mid z,a,y)h_{\theta,t,z,a}(y)
     f(F_{\theta,t}(q_{\theta,t,z},a,y))\,\sigma_a(dy)
                          \right)_{\theta,t,j,f}.
\end{equation}
Each coordinate is weak-star continuous, since its coefficient is integrable. The image is compact and convex in the product topology, including when $\Theta$ is uncountable. If the target coordinates
$p_{\theta,t+1,j}f(q_{\theta,t+1,j})$ are not in this image, a finite-coordinate projection already separates them. Indeed an open basic product neighbourhood of the target misses the compact image; in its finite projection the target is outside a compact convex set. Finite-dimensional separation then gives a linear combination of those coordinates. Absorb the finitely many scalar coefficients into affine functions $\phi_{\theta,t,j}$, using zero functions for omitted labels and including constants for masses.

The support function of \eqref{eq:uniform_moment} is exactly the right side of \eqref{eq:uniform_dual}: at each $(z,a,y)$ choose a maximizing next label \emph{after} adding all the coefficients. The least maximizing label is a measurable choice, so the integral of the maximum is attained. Thus all inequalities hold precisely when a common $k\in\mathscr K$ has every prescribed coordinate. On the standard Borel report spaces choose Borel representatives, completing the simplex on one $\sigma_a$-null set per row. Joint domination makes this completion harmless for every $\theta$, not only for a chosen countable subfamily.

Now start at $q_{\theta,0}$, for any fixed $\theta$. Assuming the proposed occupancy and conditional predictive state at cut $t$, apply the fixed action row and the common $k$. The mass coordinates in \eqref{eq:uniform_moment} give $p_{\theta,t+1,j}$. The remaining coordinates and the unnormalized instrument identity give every future-test expectation at the conditional barycenter $q_{\theta,t+1,j}$. Separating affine tests identify that barycenter. Induction realizes the whole proposal under its own deployed law. The zero-mass restrictions enforce the exact deadline and all positive-probability action restrictions. No parameter-indexed selection of a programme has occurred. Conversely, the conditional barycenters of a genuine programme obey these moment identities, and its rows are bounded by the same support function. Finally the shared terminal readouts have exactly the risks in \eqref{eq:uniform_phi}, proving \eqref{eq:uniform_exact}. The external proof keeps a separate row for each permitted phase and repeats this separation across parameters.

For example, consider one constant report under two models $\theta=0,1$, and propose terminal label $j=\theta$ with certainty. Each separate model admits that label allocation. The constant tests $\phi_{\theta,j}=\boldsymbol 1_{\{j=\theta\}}$ make the joint inequality read $2\le1$, exposing the missing common programme. Maximizing separately over labels inside the parameter sum would incorrectly accept it.

A witness is a mathematical comparison of predictions at hypothetical models. It is not an experiment permitted to query the true parameter. Lemma~\ref{lem:operational_tests} gives the relation of each affine component to executable test spans. Uniform test-span approximations within $\epsilon$ change the two-sided difference for a finite index set $\mathcal I$ by at most $2|\mathcal I|\epsilon$, since each parameter--phase occupancy has mass one. The finite witness for a \emph{proposal} is distinct from the finite global risk exclusion obtained below.

\subsection{Uniform physical presentation and proof of effectiveness}
\begin{definition}[Uniform effective presentation]\label{def:uniform_effective}
The set $\Theta$ is a computably compact metric space with supplied finite $s$-nets of uniformly computable points in $\Theta$. All experiments have a common finite-dimensional ordered physical space $(V,V_+,e)$ with compact normalized base $\mathcal B$, as in Definition~\ref{def:effective}. Preparations $x_{\theta,0}$, positive instruments $\mathsf I_{\theta,t}^a$, and losses $L_{\theta,u}$ have $0\le L_{\theta,u}\le He$. Computation of their cell integrals and losses is uniform in $\theta$. The cell partitions, reconstruction measures, support tables and terminal loss nets of Definition~\ref{def:effective} are common to the entire family. In particular, an active cell has equivalent report null sets for every normalized physical state and every parameter, and an inactive cell has zero instrument. A supplied $d_{t,h}$ bounds \eqref{eq:instrument_error} uniformly in $\theta$.

For $d_\Theta(\theta,\theta')\le s$, supplied computable moduli tending to zero give
\begin{align}\label{eq:parameter_moduli}
 \|x_{\theta,0}-x_{\theta',0}\|_{\mathcal B}&\le\omega_0(s),\\
 \sup_{a,x\in\mathcal B}\int
  \| (\mathsf I_{\theta,t}^a-\mathsf I_{\theta',t}^a)(dy)x\|_{\mathcal B}
       &\le\omega_t^{I}(s),\notag\\
 \sup_{u,x\in\mathcal B}|L_{\theta,u}(x)-L_{\theta',u}(x)|
       &\le\omega^L(s).\notag
\end{align}
Define
\begin{equation}\label{eq:parameter_Omega}
 \Omega_n(s)=\min\left\{H,\ H\left(\omega_0(s)+\sum_{t<n}\omega_t^I(s)\right)
                                  +\omega^L(s)\right\}.
\end{equation}
The errors $E_h,\beta_Q$ are those of \eqref{eq:effective_constants}, with a terminal loss net uniform in both state and parameter. Cell access, probability sampling and readout computation have the separately declared execution costs. Merely knowing that a family is continuous does not supply these certificates.
\end{definition}

\begin{theorem}[Attainment and global finite-model obstructions]\label{thm:uniform_compactness}
Suppose the raw family has a common finite-dimensional ordered physical space with compact normalized base, finite action alphabet, standard Borel reports jointly dominated by finite measures $\sigma_a$, and positive normalized instruments. Let the common readout space $U$ be compact metric, with each $u\mapsto L_{\theta,u}$ continuous on the physical space. The parameter set can be arbitrary. Every nonempty common $M$-label, exact-$n$ class has a worst-case-risk minimizer, possibly requiring a noncomputable Borel programme. Equation~\eqref{eq:finite_model_value} holds. If the common class is empty, a finite subfamily already has empty class. If $r<R^\Theta_{\aut}(n,M)$, there is a finite subfamily for which every admissible $M$-label programme has risk greater than $r$ at some one of its models. The same statements hold for fixed external widths. No bound on the size of this finite subfamily follows without additional data.
\end{theorem}
\begin{proof}
Pad every programme to $M$ labels. There are finitely many choices of its initial label and stopping subset. Fix one choice and take the product of its action simplexes, its simplex-valued $L^\infty(\sigma_a)$ report rows, and $U^M$. This is compact in the product of Euclidean, weak-star and readout topologies. A finite disjoint union over the discrete choices remains compact.

For phases at which an action is illegal, use a zero instrument only to extend the ambient occupancy formula; legality constraints exclude every execution of that extension. Thus all coefficient measures used below are dominated by the specified reference measures. Choose a basis of the finite-dimensional physical space. Every coefficient measure of $\mathsf I_{\theta,t}^a$ has an $L^1(\sigma_a)$ density. To see this, decompose a basis vector as a difference of positive vectors. Positivity bounds the total variation of its image by the sum of their masses. On a $\sigma_a$-null report set every positive image has zero mass and hence is zero, by strict positivity of the base functional; linearity handles the basis vectors. It follows that each matrix entry of
\[
 A_{\theta,t,z,a,j}(k)
       =\int k(j\mid z,a,y)\,\mathsf I_{\theta,t}^a(dy)
\]
is a weak-star continuous function of the common row. For a fixed $\theta$, propagation of the subnormalized physical-state/label vectors obeys
\[
 v_{\theta,t+1,j}=
       \sum_{z\notin D,a}\alpha_z(a)
                        A_{\theta,t,z,a,j}(k)v_{\theta,t,z}.
\]
It is a finite composition of continuous matrix and scalar operations, even when the same report row is reused. Stopping labels may be made absorbing with an arbitrary positive identity continuation solely to define these continuous functions on the entire ambient programme space. On exactly timed programmes this extension has no effect before the last cut.

At each parameter, absence of a premature stop, absence of an illegal action, and total stopping mass one at the deadline are closed equalities in these continuous occupancy functions. Intersect them over all parameters. This gives a compact admissible set. On it $R_\theta(\gamma)=\sum_j L_{\theta,u_j}(v_{\theta,n,j})$ is continuous for each fixed $\theta$. Its supremum over parameters is lower semicontinuous and thus attains its minimum on every nonempty compact admissible set. Nothing here requires the parameter to be known to the programme.

For a finite target $r$, let $C_\theta(r)$ be the closed subset of the ambient compact programme space satisfying the legality and deadline constraints at $\theta$, and $R_\theta(\gamma)\le r$. If every finite intersection of these sets is nonempty, compactness gives a point in their entire intersection. Contrapositively, a strict target below the robust optimum is excluded by a finite set of models. The same argument without the risk constraint proves finite exclusion of an empty class. Conversely the common optimum is at least every finite-subfamily optimum. Applying the just-proved finite-intersection implication to any number strictly above the supremum of those optima proves \eqref{eq:finite_model_value}, including the empty-class case. The external case has finitely many separate phase rows and the same argument. This is a global exclusion of all common programmes, not only separation of one proposed occupancy flow. It does not interchange the programme infimum and parameter supremum.
\end{proof}

\begin{lemma}[A uniform modulus over all programmes]\label{lem:uniform_modulus}
For every common admissible Borel programme $\gamma$, at every fixed persistent budget,
\begin{equation}\label{eq:uniform_modulus}
             |R_\theta(\gamma)-R_{\theta'}(\gamma)|\le\Omega_n(s)
                 \quad(d_\Theta(\theta,\theta')\le s).
\end{equation}
The constant is independent of the programme, including discontinuous report rows. It does not assume contraction of normalized predictive updates.
\end{lemma}
\begin{proof}
The physical-state/label law at a cut is a direct sum of positive subnormalized vectors $(v_z)$. Its norm is $\sum_z\|v_z\|_{\mathcal B}$. A positive mass-preserving instrument, followed by any stochastic action and report rows, contracts this norm on signed inputs: apply it to a positive decomposition and sum the preserved masses. At a fixed positive normalized input, replacing the instrument changes the next direct sum by at most $\omega_t^I(s)$, since the same report row is used and its nonnegative entries sum to one. Starting with the preparation difference, a telescoping replacement over all $n$ calls bounds the final norm by $\omega_0(s)+\sum_t\omega_t^I(s)$. The common readout has loss bounded by $H$ in dual norm; changing its physical loss functional adds at most $\omega^L(s)$. Both risks lie in $[0,H]$, proving the clipped bound. One parameter is used throughout each comparison; no timewise adversary or hidden-state trajectory for a quantum instrument is postulated.
\end{proof}

\begin{proof}[Proof of the effective part of Theorem~\ref{thm:uniform}]
For any common Borel row $k(j\mid z,a,y)$, average it against the \emph{same} reconstruction measure on each report cell. The resulting $\bar k(j\mid z,a,b)$ is independent of both phase and parameter. The positive direct-sum comparison in Lemma~\ref{lem:report_reduction}, with its error uniform in $\theta$, changes every parameter's risk by at most $E_h$. The support argument there now applies to every parameter simultaneously; therefore the one cell programme is legal and exactly timed throughout the family. Round each common row once on its original support and choose uniform-net readouts. Lemma~\ref{lem:dyadic_uniform} adds at most $\beta_Q$ to every risk and cannot create a forbidden path. Thus, if
$R_G^\Theta=\min_{\gamma\in\mathscr G_{h,Q,\rho}}\sup_\theta R_\theta(\gamma)$, then
\[
 R^\Theta_{\aut}\le R_G^\Theta\le R^\Theta_{\aut}+E_h+\beta_Q.
\]
This comparison holds for arbitrarily close Borel competitors and does not need an attained Borel optimum. The common support table proves exact emptiness equivalence with the finite list, for every supplied partition and $Q\ge1$.

Put $S_G=\min_\gamma\max_{\theta_i\in\Theta_s}R_{\theta_i}(\gamma)$. The net is a subset of the family, and Lemma~\ref{lem:uniform_modulus} is uniform in $\gamma$, so
\[
 S_G\le R_G^\Theta\le S_G+\Omega_n(s),\qquad L\le S_G\le U.
\]
For a programme $\gamma$ minimizing $\max_i l_{\gamma,i}$, every upper endpoint is at most the corresponding lower endpoint plus $\eta$. Hence $U\le L+\eta$. The risk of a minimizer of $\max_i u_{\gamma,i}$ is at most $U+\Omega_n(s)$ on the whole family. Combining the inequalities gives \eqref{eq:uniform_interval} and \eqref{eq:uniform_gap}. All individual risks are computed from the finite compositions of cell instruments, with a single $\theta_i$ in each entire composition. The candidate count is unchanged from Theorem~\ref{thm:effective}; the parameter net multiplies only the number of evaluations. Letting the four independently certified errors tend to zero proves effective convergence. At no point is $\inf_\gamma\sup_\theta$ interchanged with $\sup_\theta\inf_\gamma$.
\end{proof}

\begin{corollary}[Finite exclusion, misspecification and causal simulation]\label{cor:uniform_transfer}
For a strict target $r<R^\Theta_{\aut}(n,M)$, sufficiently fine certificates satisfy $L-E_h-\beta_Q>r$ and thus exclude \emph{every} admissible $M$-label Borel programme. The witness consists of a complete finite programme list, parameter net, risk enclosures and the uniform comparison bounds; it is not one affine continuation test.

If the true fixed parameter lies within $\delta$ of the modelled set in an ambient family satisfying the same support conditions and moduli, the attained upper bound increases by at most $\Omega_n(\delta)$. A report-implementation mismatch of total probability at most $\chi$, uniformly along the ideal deployed laws and preserving the support signature, adds at most $H\min(1,\chi)$; terminal numerical error adds its uniform loss modulus.

Suppose one parameter-independent causal simulator of a target family uses at most $S$ persistent states and $C(n)$ source acquisitions, respects the exact action/stopping/readout interface, and approximates the complete stopped and scored law within total variation $\varepsilon$, uniformly over parameters and legal target observers. Then
\begin{equation}\label{eq:uniform_simulation}
 R^\Theta_{\mathrm{source},\,\le C(n)}(MS)
       \le R^\Theta_{\mathrm{target}}(n,M)+H\varepsilon,
\end{equation}
where the source uses at most $C(n)$ calls and the source stopping interface is exactly the one in the simulator certificate. An exact-$C(n)$ source class may replace it only if that exact completion, including its clock and any padding calls, is certified. A reverse lower comparison requires a reverse simulator, not merely a one-way defect bound.
\end{corollary}
\begin{proof}
Convergence of both endpoints gives the strict exclusion. The modulus lemma compares a fixed output programme at a true parameter and a nearby modelled one. The first-discrepancy argument bounds numerical mismatch on the ideal prefix law; signature preservation is required for exact admissibility, rather than inferred from a small error probability. The simulator executes one target tuple with the product of its states and its own states. Bounded-loss total-variation comparison gives $R_\theta^{\mathrm{source}}\le R_\theta^{\mathrm{target}}+H\varepsilon$ for every $\theta$. Take the supremum and then infima over arbitrarily close target competitors. These are upper comparisons; allowed error tolerances are not automatically attained statistical lower bounds.
\end{proof}

The new minimax statement does not assert an unknown-kernel learning rate. It optimizes a parameter-independent observer over a specified, uniformly presented family, allowing whatever adaptive acquisition and learning its counted states can perform. An oracle full-history baseline depending on the true parameter cannot be subtracted as an attainable robust baseline. An excess-risk claim requires a separately specified and certified robust full-history benchmark for the same family and task. The interval theorem is for total risk and does not need such a subtraction. The singleton benchmark argument in Corollary~\ref{cor:effective_obstruction} is retained at its stated scope; it is not permission to subtract the parameterwise oracle optimum from an attained robust risk.

For the table representation, put $b=\log_2Q$ and $K_\rho=|U_\rho|$. The common programme needs at most
\[
 P(b+1)+M\lceil\log_2K_\rho\rceil+M+\lceil\log_2M\rceil
\]
bits for the probability table, readout indices, stopping flags and initial label, in addition to the declared cell and readout descriptions and fixed interpreter. The parameter net and candidate enumeration are offline costs, not retained data. Linear row sampling needs $O(b+\log(MA\max_aB_a))$ temporary bits, in addition to cell/readout scratch, and $O((A+M)(b+1))$ elementary bit operations per call in the stated read-only random-access table convention. Address formation and random access are included in that convention; a sequential-tape interpreter must instead pay its actual lookup time. Every location or random seed surviving a persistent cut is included in $Z$. These are feasible implementation bounds, not lower bounds on optimal programme size, workspace or planning time.
